% URSim fault catalogue. Generated by tools/fault_injection/catalogue.py: edit the template,
% catalogue.yaml or the captures, never the generated .tex file.
\documentclass[10pt]{article}
\usepackage[T1]{fontenc}
\usepackage[utf8]{inputenc}
\usepackage{lmodern}
\usepackage[a4paper,landscape,margin=11mm,top=13mm,bottom=13mm]{geometry}
\usepackage{microtype}
\usepackage{xcolor}
\usepackage{array}
\usepackage{longtable}
\usepackage{booktabs}
\usepackage{pifont}
\usepackage[hidelinks]{hyperref}
\usepackage{fancyhdr}

\definecolor{accent}{HTML}{4F46E5}
\definecolor{ok}{HTML}{15803D}
\definecolor{no}{HTML}{B91C1C}
\newcommand{\yes}{\textcolor{ok}{\ding{51}}}
\newcommand{\no}{\textcolor{no}{\ding{55}}}
\def\UrlBreaks{\do\/\do-\do_\do.\do=}
\urlstyle{same}
\setlength{\parindent}{0pt}
\setlength{\parskip}{4pt}
\renewcommand{\arraystretch}{1.25}
\pagestyle{fancy}\fancyhf{}
\lhead{\small URSim fault catalogue}\rhead{\small Robot Agentic}\cfoot{\small\thepage}
\renewcommand{\headrulewidth}{0.3pt}

\begin{document}

{\LARGE\bfseries\color{accent} URSim fault catalogue}\par
\vspace{2pt}
{\large Every software fault we can provoke on the Universal Robots simulator, and what it really produces}\par
\vspace{4pt}
URSim @VERSION@ (UR3e) \quad·\quad generated @DATE@ \quad·\quad @RUNS@ captures per fault

\vspace{6pt}
\textbf{Method.} Each fault is provoked by \texttt{tools/fault\_injection/inject.py} and captured by
\texttt{tools/fault\_injection/capture.py}: the robot is brought to RUNNING, our three live streams are started
(signals of the Real-Time interface at 500\,Hz, fault codes of the Primary interface, the three log files of the
container), the fault is injected, the simulator is observed for 15\,s, then the fault is reverted and the time to
RUNNING is measured. The raw evidence of every run is kept in \texttt{data/catalogue/}.

\textbf{No guessing.} Columns 3 and 4 are computed by \texttt{catalogue.py} from the captures only: each fact
gives the number of runs where it was seen (e.g.\ 5/5). Log lines are shown as parsed by the agent's own
connector, with numbers replaced by \#. The official meaning of a fault code comes from the Universal Robots error
code pages (software 5.26). Columns 1, 2 and 5 are written by hand in \texttt{catalogue.yaml}: commands copied
from the injection script, and official sources, each with what it supports. Facts observed outside the dataset are
marked as notes, with how they were observed.

\vspace{4pt}
\begin{tabular}{@{}llcccc@{}}
\toprule
\textbf{Fault} & \textbf{Category} & \textbf{Fault code} & \textbf{Warning/error in a log} & \textbf{Visible in the signals} & \textbf{Black box} \\
\midrule
@OVERVIEW@
\bottomrule
\end{tabular}

\smallskip
{\small ``Visible in the signals'': a stream cut, a change of robot/safety mode, or a median rate below 400\,Hz
(500\,Hz by design). A fault with only \no\ in the first two columns cannot raise an alert from logs and fault codes
alone: it needs the rules on the signals (summary.md, section~3).}

\newpage
\footnotesize
\begin{longtable}{@{}>{\raggedright\arraybackslash}p{2.7cm}>{\raggedright\arraybackslash}p{4.0cm}>{\raggedright\arraybackslash}p{7.6cm}>{\raggedright\arraybackslash}p{5.2cm}>{\raggedright\arraybackslash}p{6.0cm}@{}}
\toprule
\textbf{Fault} & \textbf{How to provoke / revert} & \textbf{Logs and fault codes (observed)} & \textbf{Effect on signals (observed)} & \textbf{Category, sources and proof} \\
\midrule
\endhead
@ROWS@
\end{longtable}

\end{document}
